\chapter{Bonus Proofs}

\section{\texorpdfstring{$\text{ZFC}+\text{MA}\vdash\mf{dd}=\mf{c}$}{MA=>dd=c}}

We already knew that MA implies $\mf{dd}=\mf{c}$, since, as we saw, $\text{non}(\mc{M})\leq\mf{dd}$, and we know MA collapses Cichoń's Diagram. However, let's do a direct proof of it by exhibiting a suitable poset.

\begin{theorem}{}{MA_dd=c}
    $\text{ZFC}+\text{MA}(\kappa)\vdash\mf{dd}>\kappa$.
\end{theorem}

So, given a set $\mc{F}\subseteq\sym$, with $\lvert\mc{F}\rvert\leq\kappa$, using $\text{MA}(\kappa)$ we will find an infinite and coinfinite subset $A\subseteq\omega$ such that $d(\pi[A])=d(A)=0$, for all $\pi\in\mc{F}$. In order to simplify our work, let's assume WLOG that $\text{id}\in\mc{F}$.

Intuitively, the idea to construct such set $A$ is to partition $\omega$ as
$$\omega=\bigsqcup_{k\in\omega}I_k,\text{ where }\lvert I_k\rvert/\lvert I_{k+1}\rvert\to0$$
and let $A$ be a set such that $\pi[A]$ has at most one point in each interval $I_k$.

Throughout the proof, we will use interchangeably the characteristic function $s_A\in2^{<\omega}$ of our generic set $A\subseteq\omega$ we want to construct, and the set itself.

\begin{sketch}
    Before proving it directly, let's construct our poset intuitively. We want our set of conditions $\mbb{P}$ to be such that $p=(s_p,F_p,n_p)$, where $s_p\in2^{<\omega}$, $F_p\in[\mc{F}]^{<\omega}$ and $n_p\in\omega$. And maybe the most naive common way to define $p\leq q$ is as $s_p\supseteq s_q$,$F_p\supseteq F_q$,$n_p\geq n_q$, and
    $$\lvert\pi[A_p]\cap I_k\rvert\leq 1,\;\forall k\geq n_q,\;\forall\pi\in F_q$$
    Let's see the possible problems that could arise with such definition. Clearly $(\mbb{P},\leq)$ is antisymmetric, but what about reflexivity? In fact, nothing guarantees us that $\lvert\pi[A_p]\cap I_k\rvert\leq 1$ holds for all $\pi\in F_p$ and $k\geq n_p$, so what we need to do is to add such condition in the definition of $\mbb{P}$. Having done that, let's try to prove transitivity.

    Given $p\leq q\leq r$, let $\pi\in F_r$ and $k\geq n_r$. If $k\geq n_q$, this automatically follows from $p\leq q$, otherwise, $n_r\leq k<n_q$. In this case, by $q\leq r$ we have that
    $$\lvert\pi[A_q]\cap I_k\rvert\leq 1$$
    but we want $\pi[A_p]$ there instead of $\pi[A_q]$. In order to do that, we will state in the definition of $\leq$ that every new element in the extension isn't in any $I_k$, where $k<n_q$, which will clearly be sufficient to guarantee that $\pi[A_q]\cap I_k=\pi[A_p]\cap I_k$ for those $k$'s.
\end{sketch}

\begin{proof}
    Let
    $$\mbb{P}=\left\{p=(s_p,F_p,n_p):s_p\in2^{<\omega},F_p\in[\mc{F}]^{<\omega},n_p\in\omega,\;\forall\pi\in F_p\;\forall k\geq n_p\;(\lvert\pi[A_p]\cap I_k\rvert\leq 1\right)\}$$
    And consider $(\mbb{P},\leq)$, where $p\leq q$ iff $s_p\supseteq s_q$, $F_p\supseteq F_q$, $n_p\geq n_q$ and, for all $\pi\in F_q$:
    \begin{enumerate}
        \item $\lvert\pi[A_p]\cap I_k\rvert\leq 1$, $\forall k\geq n_q$;
        \item $\pi(n)\notin\bigcup_{k<n_q}I_k$, $\forall n\in A_p\setminus A_q$.
    \end{enumerate}
    Clearly $\leq$ is antisymmetric and reflexive, since (1) is guaranteed by definition of $\mbb{P}$, and (2) is trivial for $A_p=A_q$. For transitivity, (1) was proved in the previous sketch, since condition (2) is added so (1) could hold, then it's sufficient to show (2). In fact, if $p\leq q\leq r$, given $n\in A_p\setminus A_r$ we have two options:
    \begin{itemize}
        \item If $n\in A_q\setminus A_r$, then it follows imediatelly from $q\leq r$;
        \item If $n\in A_p\setminus A_q$, then $\pi(n)\notin\bigcup_{k<n_q}I_k\supseteq\bigcup_{k<n_r}I_k$.
    \end{itemize}

    Now, $(\mbb{P},\leq)$ is clearly $\sigma$-centered, since
    $$\mbb{P}=\bigcup_{n\in\omega,\,s\in2^{<\omega}}\left\{(s,F,n)\in\mbb{P}:F\in[\mc{F}]^{<\omega}\right\}$$
    and $(s,F_1,n)\leq(s,F_2,n)$ iff $F_1\supseteq F_2$. So $(\mbb{P},\leq)$ is ccc, as desired.

    Consider now the sets, for $m\in\omega$ and $\sigma\in\mc{F}$:
    $$D_m=\{p\in\mbb{P}:\lvert A_p\rvert\geq n\}\qquad\text{and}\qquad D_\sigma=\{p\in\mbb{P}:\sigma\in F_p\}$$
    If they were dense, then we can apply $MA(\kappa)$ for $\mc{D}=\{D_m,D_\sigma:m\in\omega,\;\sigma\in\mc{F}\}$, since $\lvert\mc{D}\rvert\leq \kappa$. Fix then $q\in\mbb{P}$, and consider $m\in\omega$ and $\sigma\in\mc{F}$, then:
    \begin{itemize}
        \item To see $D_m$ is dense, it's sufficient to add one element $n$ to $A_q$, i.e., find $p\leq q$ such that $F_p=F_q$, $n_p=n_q$ and $A_p=A_q\cup\{n\}$, where $n\notin A_q$. In order to do that, let's see which $n$ we need such that $p\leq q$. We need that, for every $\pi\in F_q$, that
        \begin{itemize}
            \item $\pi[A_p\setminus A_q]\cap\bigcup_{k<n_q}I_k=\emptyset$. Then it's sufficient to show $\pi(n)\notin\bigcup_{k<n_q}I_k$;
            \item $\lvert\pi[A_q]\cap I_k\rvert\leq 1$, $\forall k\geq n_p$, i.e., if $\pi[A_q]\cap I_k\neq\emptyset$, then $\pi(n)\notin I_k$.
        \end{itemize}
        But notice that, since we have a finite amount of $\pi\in F_q$, and both $\bigcup_{k<n_q}I_k$ and the amount of intervals $I_k$ where $\pi[A_q]\cap I_k\neq\emptyset$ are finite, we can let $n$ be any number greater than all of those.
        \item To see $D_\sigma$ is dense, let $p=(s_q, F_q \cup\{\sigma\}, n_p)\leq q$, where $n_p\geq n_q$ is such that
        $$\sigma[A_p]=\sigma[A_q]\subseteq\bigcup_{k<n_p}I_k$$
        which implies that $\lvert\sigma[A_p]\cap I_k\rvert=0\leq 1$, $\forall k\geq n_p$, so $p\in\mbb{P}$. It's not hard to see we can choose such $n_p$, since $\sigma[A_q]$ is finite. Then we have $p\leq q$, as desired.
    \end{itemize}
    Now, since $\mc{D}$ is a family of $\leq\kappa$ dense sets, by $MA(\kappa)$ we can find a $\mc{D}$-generic filter $G$ in $\mbb{P}$. Let
    $$A_G=\bigcup_{p\in G}A_p$$
    Clearly, since $G\cap D_n\neq\emptyset$ for every $n\in\omega$, then $A_G$ is infinite, so it's a non-trivial object. Now, given $\pi\in\mc{F}$, let $q\in G\cap D_\pi$, we claim that $\lvert\pi[A_G]\cap I_k\rvert\leq 1$, for every $k\geq n_q$.

    In fact, let $p$ be an element extending both $q$ and an element in $G\cap D_n$, with $n>k$ large enough such that $\pi^{-1}[I_k]\subseteq\text{dom}(s_p)$, this guarantees us that $\pi[A_G]\cap I_k=\pi[A_p]\cap I_k$, since $s_p\subseteq s_G$. But then, since we have $p\leq q$, it implies that
    $$\lvert\pi[A_G]\cap I_k\rvert=\lvert\pi[A_p]\cap I_k\rvert\leq 1$$
    as desired, so $d(\pi[A_G])=0$.
\end{proof}